\documentclass[a4paper]{article}
\usepackage[margin=1in]{geometry}
\usepackage{amsmath,amsfonts,mathtools,amsthm,amssymb}
\usepackage[l2tabu,orthodox]{nag}
\usepackage{microtype} % fixes spacing or whatever
\usepackage{enumitem}
\usepackage{tikz-cd}
\usepackage{xcolor}
\usepackage{physics}
\usepackage{mathrsfs}
\usepackage{cancel}
% \usepackage{libertine,libertinust1math}
\usepackage{eucal}
\usepackage[toc,page]{appendix}
\usepackage{pgfplots}
\pgfplotsset{compat=1.18}

\newtheorem{proposition}{Proposition}
\newtheorem{theorem}{Theorem}
\newtheorem{lemma}{Lemma}
\newtheorem{corollary}{Corollary}

\theoremstyle{definition}
\newtheorem{remark}{Remark}
\newtheorem{definition}{Definition}
\newtheorem{example}{Example}

% \usepackage[sc,osf]{mathpazo}   % With old-style figures and real smallcaps.
% \linespread{1.025}              % Palatino leads a little more leading
% % Euler for math and numbers
% \usepackage[euler-digits,small]{eulervm}
% \AtBeginDocument{\renewcommand{\hbar}{\hslash}}

\usepackage{etoolbox}
\usepackage{dynkin-diagrams}
\def\row#1/#2!{\dynkin{#1}{#2}\\}
\newcommand{\tble}[1]{
   \renewcommand*\do[1]{\row##1!}
   \[
      \begin{array}{ll}\docsvlist{#1}\end{array}
   \]
}

\allowdisplaybreaks

\renewcommand{\epsilon}{\varepsilon}
% \renewcommand{\phi}{\varphi}

\DeclareMathAlphabet{\mathpzc}{OT1}{pzc}{m}{it}

\newcommand{\goth}[1]{\mathfrak{#1}}
\newcommand{\red}[1]{#1_{\text{red}}}
\newcommand{\ad}[1]{#1_{\text{Ad}}}
\newcommand{\reg}[1]{#1_{\text{reg}}}
\newcommand{\sh}[1]{#1^{\text{sh}}}
\newcommand{\cpdv}[2]{\cfrac{\partial #1}{\partial #2}}

\newcommand{\fA}{\mathcal{A}}
\newcommand{\fB}{\mathcal{B}}
\newcommand{\fC}{\mathcal{C}}
\newcommand{\fD}{\mathcal{D}}
\newcommand{\fE}{\mathcal{E}}
\newcommand{\fF}{\mathcal{F}}
\newcommand{\fG}{\mathcal{G}}
\newcommand{\fH}{\mathcal{H}}
\newcommand{\fI}{\mathcal{I}}
\newcommand{\fJ}{\mathcal{J}}
\newcommand{\fK}{\mathcal{K}}
\newcommand{\fL}{\mathcal{L}}
\newcommand{\fM}{\mathcal{M}}
\newcommand{\fN}{\mathcal{N}}
\newcommand{\fO}{\mathcal{O}}
\newcommand{\fP}{\mathcal{P}}
\newcommand{\fQ}{\mathcal{Q}}
\newcommand{\fR}{\mathcal{R}}
\newcommand{\fS}{\mathcal{S}}
\newcommand{\fT}{\mathcal{T}}
\newcommand{\fX}{\mathcal{X}}
\newcommand{\fY}{\mathcal{Y}}
\newcommand{\fZ}{\mathcal{Z}}

\newcommand{\be}{\mathbf{e}}
\newcommand{\bx}{\mathbf{x}}
\newcommand{\bone}{\mathbf{1}}
\newcommand{\cx}{\bar{x}}
\newcommand{\cy}{\bar{y}}
\newcommand{\vx}{\vec{x}}
\newcommand{\vsigma}{\vec{\sigma}}
\newcommand{\tzeta}{\tilde{\zeta}}

\newcommand{\gm}{\goth{m}}
\newcommand{\gp}{\goth{p}}
\newcommand{\gF}{\goth{F}}
\newcommand{\gU}{\goth{U}}
\newcommand{\gV}{\goth{V}}

\newcommand{\A}{\mathbf{A}}
\newcommand{\C}{\mathbf{C}}
\newcommand{\F}{\mathbf{F}}
\newcommand{\G}{\mathbf{G}}
\newcommand{\bH}{\mathbf{H}}
\newcommand{\N}{\mathbf{N}}
\newcommand{\PP}{\mathbf{P}}
\newcommand{\Q}{\mathbf{Q}}
\newcommand{\R}{\mathbf{R}}
\newcommand{\Z}{\mathbf{Z}}

\newcommand{\dlog}{\dd\log}

\newcommand\srestr[2]{{\left.\kern-\nulldelimiterspace #1\vphantom{\small|} \right|_{#2}}}
\newcommand\restr[2]{{\left.\kern-\nulldelimiterspace #1 \vphantom{\big|} \right|_{#2}}}
\newcommand\gsec[2]{\Gamma({#1}, {#2})}
\newcommand\gsecx[1]{\Gamma(X, {#1})}

\DeclareMathOperator{\chr}{char}
\DeclareMathOperator{\rProj}{\mathbf{Proj}}
\DeclareMathOperator{\rSpec}{\mathbf{Spec}}
\DeclareMathOperator{\rHom}{\mathbf{Hom}}
\DeclareMathOperator{\rExt}{\mathbf{Ext}}
\DeclareMathOperator{\id}{id}
\DeclareMathOperator{\Frac}{Frac}
\DeclareMathOperator{\rk}{rank}
\DeclareMathOperator{\deter}{det}
\DeclareMathOperator{\pic}{Pic}
\DeclareMathOperator{\cacl}{CaCl}
\DeclareMathOperator{\trd}{tr.d.}
\DeclareMathOperator{\cl}{Cl}
\DeclareMathOperator{\depth}{depth}
\DeclareMathOperator{\codim}{codim}
\DeclareMathOperator{\Div}{Div}
\DeclareMathOperator{\coker}{coker}
\DeclareMathOperator{\len}{length}
\DeclareMathOperator{\height}{height}
\DeclareMathOperator{\supp}{Supp}
\DeclareMathOperator{\proj}{Proj}
\DeclareMathOperator{\im}{im}
\DeclareMathOperator{\Hom}{Hom}
\DeclareMathOperator{\Der}{Der}
\DeclareMathOperator{\spec}{Spec}
\DeclareMathOperator{\Aut}{Aut}
\DeclareMathOperator{\ch}{char}
\DeclareMathOperator{\tor}{Tor}
\DeclareMathOperator{\Ann}{Ann}
\DeclareMathOperator{\Syl}{Syl}
\DeclareMathOperator{\Sym}{Sym}
\DeclareMathOperator{\GL}{GL}
\DeclareMathOperator{\SL}{SL}
\DeclareMathOperator{\Stab}{Stab}
\DeclareMathOperator{\Perm}{Perm}
\DeclareMathOperator{\Orb}{Orb}
\DeclareMathOperator{\Gal}{Gal}
\DeclareMathOperator{\Supp}{Supp}
\DeclareMathOperator{\Ext}{Ext}
\DeclareMathOperator{\Tor}{Tor}
\DeclareMathOperator{\PGL}{PGL}
\DeclareMathOperator{\hd}{hd}
\DeclareMathOperator{\pd}{pd}
\DeclareMathOperator{\SO}{SO}
\DeclareMathOperator{\SU}{SU}
\DeclareMathOperator{\Sp}{Sp}
\DeclareMathOperator{\Oth}{O}
\DeclareMathOperator{\Uni}{U}
\DeclareMathOperator{\evf}{ev}
\DeclareMathOperator{\cH}{H}
\DeclareMathOperator{\dR}{R}
\DeclareMathOperator{\End}{End}
\DeclareMathOperator{\Hasse}{Hasse}
\DeclareMathOperator{\Num}{Num}

\author{James Lee}

% Solarized
% \pagecolor[RGB]{0,20,26}
% \color[RGB]{191,191,191}

% Sephia
% \pagecolor[RGB]{249,239,220}

% Grey
% \pagecolor[RGB]{200, 200, 200}

% Black
% \pagecolor[RGB]{0,0,0}
% \color[RGB]{255,255,255}

\title{Chapter 3, Section 8}

\begin{document}
\maketitle
\begin{enumerate} [label=\textbf{\arabic*.}, leftmargin=0em]

  \item Let $f : X \to Y$ be a continuous map of topological spaces.
        Let $\fF$ be a sheaf of Abelian groups on $X$, and assume that $\dR^if_*(\fF) = 0$ for all $i > 0$.
        Show that there are natural isomorphisms, for each $i \geq 0$,
        \begin{equation*}
          \cH^i(X, \fF) \cong \cH^i(Y, f_* \fF).
        \end{equation*}

        \begin{proof}
          Let $0 \to \fF \to \fI^\bullet$ be an injective resolution of $\fF$.
          By hypothesis, $0 \to f_* \fF \to f_* \fI^\bullet$ is also exact.
          The pushforward of an injective object in $\goth{Ab}(X)$ is also injective in $\goth{Ab}(Y)$ because $f_*$ is a right adjoint functor, so $0 \to f_*\fF \to f_* \fI^\bullet$ is an injective resolution of $f_* \fF$ in $\goth{Ab}(Y)$.
          Suppose $\fG \to \fG'$ is a monomorphism of sheaves on $Y$, and $\fG \to f_*\fI$ is any morphism of sheaves.
          By adjunction, there is a morphism $f^{-1} \fG \to \fI$, and because $f^{-1}$ is an exact functor, $f^{-1} \fG \to f^{-1} \fG'$ is also a monomorphism (note that in general, $f^*$ is only right exact).
          By injectivity of $\fI$, there is an extension of $f^{-1} \fG \to \fI$ making the following diagram commute
          \[ \begin{tikzcd}
              & \fI                           \\
              f^{-1} \fG \arrow[r] \arrow[ru] & f^{-1} \fG' \arrow[u, dashed]
            \end{tikzcd}. \]
          Hence, by adjunction again, there is an extension of $\fG \to f_* \fI$ making the following diagram commute
          \[ \begin{tikzcd}
              & f_*\fI                           \\
              \fG \arrow[r] \arrow[ru] & \fG' \arrow[u, dashed]
            \end{tikzcd}. \]
          Since $\fF$ and $f_*\fF$ have same global sections for any $\fF$, there are natural identificiations
          \[ \begin{tikzcd}
              0 \arrow[r] & {\Gamma(X, \fF)} \arrow[r] \arrow[d, Rightarrow, no head] & {\Gamma(X, f_* \fI^\bullet)} \arrow[d, Rightarrow, no head] \\
              0 \arrow[r] & {\Gamma(Y, f_* \fF)} \arrow[r]                            & {\Gamma(Y, f_*\fI^\bullet)}
            \end{tikzcd}. \]
          Taking cohomology gives the desired natural isomorphisms for all $i \geq 0$.
        \end{proof}

  \item Let $f : X \to Y$ be an affine morphism of schemes (II, Ex. 5.17) with $X$ Noetherian, and let $\fF$ be a quasi-coherent sheaf on $X$.
        Show that the hypotheses of (Ex. 8.1) are satisfies, and hence that $\cH^i(X, \fF) \cong \cH^i(Y, f_* \fF)$ for each $i \geq 0$.

        \begin{proof}
          Let $V$ be an open affine subset of $Y$.
          By (8.2) and (8.5), we have
          \begin{equation*}
            \restr{\dR^i f_*(\fF)}{V} \cong \widetilde{\cH^i(f^{-1}V, \restr{\fF}{f^{-1}(V)})},
          \end{equation*}
          which is zero because $f^{-1}(V)$ is affine (3.7).
          Hence, $\dR^i f_*(\fF) = 0$, and the result follows by the previous exercise.
        \end{proof}

  \item Let $f : X \to Y$ be a morphism of ringed spaces, let $\fF$ be an $\fO_X$-module, and let $\fE$ be a locally free $\fO_Y$-module of finite rank.
        Prove the \textit{projection formula} (cf. (II, Ex. 5.1))
        \begin{equation*}
          \dR^if_*(\fF \otimes f^* \fE) \cong \dR^i f_*(\fF) \otimes \fE.
        \end{equation*}

        \begin{proof}
          The pullback of a locally free sheaf of finite rank is also locally free of finite rank.
          Localization is an exact monoidal functor from the category of $\fO_Y$-modules to the category of Abelian groups, so by checking exactness at stalks, tensoring by $\fE$ is an exact functor, and it also has a left adjoint, where for any $\fO_Y$-modules $\fM$ and $\fN$ (II, Ex. 5.1)
          \begin{align*}
            \Hom_{\fO_Y}(\fM \otimes \fE^\vee, \fN) & \cong \Hom_{\fO_Y}(\fM, \rHom_{\fO_Y}(\fE^\vee, \fN)) \\
                                                    & \cong \Hom_{\fO_Y}(\fM, \fN \otimes (\fE^\vee)^\vee)  \\
                                                    & \cong \Hom_{\fO_Y}(\fM, \fN \otimes \fE).
          \end{align*}
          Hence, the result follows from the $i = 0$ case, (1.3A), and Lemma 1.

          It remains to show
          \begin{equation*}
            f_*(\fF \otimes f^* \fE) \cong f_*(\fF) \otimes \fE.
          \end{equation*}
          We first show there is a natural map
          \begin{equation*}
            \phi : f_*(\fF) \otimes \fE \to f_*(\fF \otimes f^*\fE).
          \end{equation*}
          Let $\fP$ be the presheaf on $Y$ defined for open subsets $V \subseteq Y$ by
          \begin{equation*}
            V \mapsto f_*(\fF)(V) \otimes \fE(V)
          \end{equation*}
          so that $\sh{\fP} \cong f_*(\fF) \otimes \fE$, and let $\fQ$ be the presheaf on $X$ defined for open subsets $U \subseteq X$ by
          \begin{equation*}
            U \mapsto \fF(U) \otimes f^*\fE(U)
          \end{equation*}
          so that $\sh{\fQ} \cong \fF \otimes f^*\fE$.
          Sheafification and pushforward do not commute; nonetheless there is a natural morphism $\sh{(f_*\fQ)} \to f_*(\sh{\fQ})$.
          In particular, consider the natural morphism $\fQ \to \sh{\fQ}$.
          Taking the pushforward gives $f_*\fQ \to f_*(\sh{\fQ})$.
          Because $f_*(\sh{\fQ})$ is already a sheaf, by the universal property of sheafification (II, 1.2), there is unique map $\tau : \sh{(f_*\fQ)} \to f_*(\sh{\fQ})$ fitting into the commutative diagram
          \[ \begin{tikzcd}
              & \sh{(f_*\fQ)} \arrow["\tau", rd, dashed] &               \\
              f_*\fQ \arrow[rr] \arrow[ru] &                                & f_*(\sh{\fQ})
            \end{tikzcd}. \]
          Thus, we define a morphism of presheaves $\psi : \fP \to f_*\fQ$, which induces a morphism of sheaves $\sh{\psi} : \sh{\fP} \to \sh{(f_*\fQ)}$, and composing with $\tau$ gives our desired morphism
          \begin{equation*}
            \phi = \tau \circ \sh{\psi} : \sh{\fP} \to \sh{(f_*\fQ)} \to f_*(\sh{\fQ}).
          \end{equation*}
          The pushforward for the presheaf $\fQ$ distributes over the tensor product, that is
          \begin{equation*}
            \Gamma(V, f_*\fQ) = f_*(\fF)(V) \otimes f_*f^*(\fE)(V).
          \end{equation*}
          There is a natural morphism of sheaves $\iota : \fE \to f_*f^*(\fE)$, so take $\psi = \id_{f_*\fF} \otimes \iota$ to be our morphism of presheaves
          \[ \begin{tikzcd}
              \sh{\fP} \arrow[r, "\sh{\psi}"] \arrow[rr, "\phi", bend left] & \sh{(f_*\fQ)} \arrow[r, "\tau"] & f_*(\sh{\fQ}) \\
              \fP \arrow[r, "\psi"'] \arrow[u]                              & f_*\fQ \arrow[u] \arrow[ru]     &
            \end{tikzcd}. \]
          Now that we defined a morphism, we remark that the question is local on $Y$.
          In particular, let $y$ be a point in $Y$, and let $V$ be an open subset of $Y$ such that $\fE$ is free of finite rank.
          Then $f^*\fE$ is free of finite rank on $U = f^{-1}(V)$.
          Consider the commutative diagram
          \[ \begin{tikzcd}
              U \arrow[r, "\bar{f}"] \arrow[d, "i"'] & V \arrow[d, "j"] \\
              X \arrow[r, "f"]                       & Y
            \end{tikzcd}. \]
          Because $i$ and $j$ are open immersions, there is an isomorphism of functors
          \begin{equation*}
            j^*f_* \cong \bar{f}_*i^*, \quad i^*f^* \cong \bar{f}^*j^*.
          \end{equation*}
          Pullback distributes over tensor products, so we have
          \begin{align*}
            j^*f_*(\fF \otimes f^*\fE) & \cong \bar{f}_*i^*(\fF \otimes f^*\fE)           \\
                                       & \cong \bar{f}_*(i^*\fF \otimes i^*f^*\fE)        \\
                                       & \cong \bar{f}_*(i^*\fF \otimes \bar{f}^*j^*\fE), \\
            j^*(f_*(\fF) \otimes \fE)  & \cong j^*f_*(\fF) \otimes j^*\fE                 \\
                                       & \cong \bar{f}_*(i^*\fF) \otimes j^*\fE.
          \end{align*}
          Open immersions preserve stalks, so we can assume $\fE$ to be a free $\fO_X$-module of finite rank $n$.
          The pullback of a free sheaf of finite rank is also free of finite rank, so we have the isomorphisms
          \begin{align*}
            f_*(\fF \otimes f^*\fE) & \cong f_*(\fF \otimes \fO_X^{\oplus n}) \\
                                    & \cong f_*(\fF^{\oplus n})               \\
                                    & \cong f_*(\fF)^{\oplus n}               \\
                                    & \cong f_*(\fF) \otimes \fO_Y^{\oplus n} \\
                                    & \cong f_*(\fF) \otimes \fE.
          \end{align*}
          The third isomorphism holds because $f_*$ is right adjoint, and right adjoint functors preserve limits, and a direct sum of finite number of terms can be realized as a limit, so the pushforward and finite direct sums commute.
        \end{proof}

        \begin{lemma}
          Let $\fA, \fB, \fC$ be Abelian categories, and assume $\fA$ and $\fB$ have enough injectives.
          Let $F : \fA \to \fB$ and $G : \fB \to \fC$ be functors such that $F$ has a left adjoint (or more generally preserves injectives), and either
          \begin{itemize}
            \item[(a)] $F$ is left exact, and $G$ is exact; or
            \item[(b)] $F$ is exact, and $G$ is left exact.
          \end{itemize}
          In each case, we have the following natural isomorphisms
          \begin{itemize}
            \item[(a)] $\dR^i(G \circ F) \cong G \circ \dR^i(F)$;
            \item[(b)] $\dR^i(G \circ F) \cong \dR^i(G) \circ F$.
          \end{itemize}
        \end{lemma}

        \begin{proof}
          The proof for (b) is similar to the proof of Exercise 1, so we prove (a).
          Let $A$ be an object in $\fA$ and choose an injective resolution
          \begin{equation*}
            0 \to A \to I^\bullet,
          \end{equation*}
          where one has
          \begin{equation*}
            \dR^iF(A) = h^i(FI^\bullet).
          \end{equation*}
          Applying $G \circ F$ gives a complex in $\fC$
          \begin{equation*}
            0 \to (G \circ F)A \to (G \circ F) I^\bullet,
          \end{equation*}
          and we want to show
          \begin{equation*}
            h^i((G \circ F)I^\bullet) \cong G \circ h^i(FI^\bullet).
          \end{equation*}
          But this follows from exact functors commuting with cohomology.

        \end{proof}

\end{enumerate}

\end{document}

% \item Let $Y$ be a Noetherian scheme, and let $\mathscr{E}$ be a locally free $\fO_Y$-module of rank $n + 1$, $n \geq 1$. Let $X = \PP(\mathscr{E})$ (II, \S 7), with the invertible sheaf $\fO_X(1)$ and the projection morphism $\pi : X \to Y$.
% \begin{itemize}
%   \item[(a)] Then $\pi_*(\fO(l)) \cong S^l(\fE)$ for $l \geq 0$, $\pi_*(\fO(l)) = 0$ for $l < 0$ (II, 7.11); $R^i\pi_*(\fO(l)) = 0$ for $0 < i < n$ and all $l \in \Z$; and $R^n \pi_*(\fO(l)) = 0$ for $l > - n - 1$.
%   \item[(b)] Show there is a natural exact sequence
%   \begin{equation*}
%     0 \to \Omega_{X / Y} \to (\pi^* \fE)(-1) \to \fO \to 0,
%   \end{equation*}
%   cf. (II, 8.13), and conclude that the \textit{relative canonical sheaf} $\omega_{X/Y} = \bigwedge^n \Omega_{X / Y}$ is isomorphic to $(\pi^* \bigwedge^{n + 1} \fE)(-n - 1)$. Show furthermore that there is a natural isomorphism $R^n \pi_*(\omega_{X / Y}) \cong \fO_Y$ (cf. (7.1.1)).
%   \item[(c)] Now show, for any $l \in \Z$, that
%   \begin{equation*}
%     R^n \pi_*(\fO(l)) \cong \pi_*(\check{\fO(-l - n - 1)}) \otimes \check{\bigwedge^{n + 1} \fE}
%   \end{equation*}
%   \item[(d)] Show that $p_a(X) = (-1)^n p_a(Y)$ (use (Ex. 8.1)) and $p_g(X) = 0$ (use II, 8.11).
%   \item[(e)] In particular, if $Y$ is a nonsingular projective curve of genus $g$, and $\fE$ a locally free sheaf of rank $2$, then $X$ is a projective surface with $p_a = -g$, $p_g = 0$, and irregularity $g$ (7.12.3). This kind of surface is called a \textit{geometrically ruled surface} (V, \S 2).
% \end{itemize}
